% --- Tabel Deskripsi UC-06 ---
\begin{table}[H]
    \centering
    \caption{Deskripsi Use Case Memantau Konteks Sosial (UC-06)}
    \label{tbl:desc_uc06}
    \begin{tabular}{|p{0.20\textwidth}|p{0.70\textwidth}|}
        \hline
        \textbf{Nama Use Case} & \textbf{Memantau Konteks Sosial} \\
        \hline
        \textbf{Aktor} & pelajar \\
        \hline
        \textbf{FR Terkait} & FR-04 \\
        \hline
        \textbf{Deskripsi} & pelajar melihat posisi relatif kemahirannya dibandingkan dengan distribusi kemampuan kelas secara anonim (\textit{Social Awareness}). \\
        \hline
        \textbf{Kondisi sebelum} & Data agregat kemampuan seluruh kelas tersedia. \\
        \hline
        \textbf{Kondisi sesudah} & Visualisasi distribusi kelas ditampilkan. \\
        \hline
        \textbf{Main Flow} & \begin{enumerate}[noitemsep,topsep=0pt,leftmargin=*]
            \item pelajar mengakses menu wawasan grup/kelas.
            \item Sistem mengagregasi data pengguna aktif.
            \item Sistem menampilkan distribusi nilai grup/kelas.
        \end{enumerate} \\
        \hline
    \end{tabular}
\end{table}